\renewcommand{\arraystretch}{1.15}
\begin{longtable}{@{}p{0.28\textwidth}p{0.68\textwidth}@{}}
\toprule
\textbf{Terme} & \textbf{Définition} \\
\midrule
\endhead

\textbf{App Router (Next.js)} & Système de routage de Next.js basé sur le dossier \texttt{app/}, utilisant des \textit{Server Components} par défaut et des \textit{Route Handlers} pour l’API. \\

\textbf{API (REST)} & Interface permettant à un client (navigateur) d’échanger avec le serveur via des requêtes HTTP (\texttt{GET}, \texttt{POST}, \texttt{PATCH}, \texttt{DELETE}) et des réponses structurées (souvent en JSON). \\

\textbf{Back-office / Admin} & Partie de l’application réservée à l’administration : gestion des produits, catégories, services, rendez-vous, messages, blog et médias. \\

\textbf{Cloudinary} & Service de stockage et de distribution d’images/vidéos via CDN, avec possibilités d’optimisation et de transformation à la volée. \\

\textbf{CDN} & \textit{Content Delivery Network} : réseau de serveurs distribué géographiquement qui met en cache et délivre des contenus (ex. images) au plus près des utilisateurs, afin de réduire la latence et d’améliorer les performances. \\

\textbf{CRUD} & Ensemble d’opérations standards sur des données : \textit{Create, Read, Update, Delete}. \\

\textbf{ERD} & \textit{Entity Relationship Diagram} : diagramme entité-association représentant les tables, attributs et relations d’une base de données relationnelle. \\

\textbf{SEO} & \textit{Search Engine Optimization} : ensemble de pratiques visant à améliorer la visibilité d’un site dans les moteurs de recherche (structure des pages, performance, contenus). \\

\textbf{SGBD} & \textit{Système de Gestion de Base de Données} : logiciel permettant de stocker, organiser et interroger des données (ex. PostgreSQL). \\

\textbf{Middleware} & Mécanisme exécuté avant le traitement principal d’une requête (ex. ajout d’en-têtes de sécurité, protection d’accès), pouvant s’appliquer à des routes spécifiques. \\

\textbf{Neon} & Fournisseur de PostgreSQL managé (compatible serverless), souvent utilisé avec des plateformes de déploiement comme Vercel. \\

\textbf{ORM (Prisma)} & Couche d’abstraction entre la base SQL et le code applicatif, permettant d’écrire des requêtes via un client typé plutôt qu’en SQL brut. \\

\textbf{PostgreSQL} & SGBD relationnel robuste supportant contraintes, transactions, index, et adapté à des données structurées (rendez-vous, produits, articles). \\

\textbf{SSR} & \textit{Server-Side Rendering} : génération HTML côté serveur pour améliorer SEO et performance perçue. \\

\textbf{TypeScript} & Sur-ensemble de JavaScript ajoutant un typage statique. Les types sont vérifiés à la compilation, puis effacés à l’exécution. \\

\textbf{Vercel} & Plateforme de déploiement optimisée pour Next.js, permettant CI/CD, prévisualisations, et hébergement global. \\

\textbf{CI/CD} & \textit{Continuous Integration / Continuous Delivery} : automatisation des étapes de test, build et déploiement à chaque mise à jour du code. \\

\bottomrule
\end{longtable}
\renewcommand{\arraystretch}{1.0}
